%%%%%%%%%%%%%%%%%%%%%%%%%%%%%%%%%%%
% Supporting Information — Digital Discovery
% Repositioned 2026-06-23: honest structure-based benchmark (see HONEST_ML_ASSESSMENT_260622.md)
%%%%%%%%%%%%%%%%%%%%%%%%%%%%%%%%%%%

\documentclass[twoside,twocolumn,9pt]{extarticle}
\usepackage{extsizes}
\usepackage[super,sort&compress,comma]{natbib}
\usepackage[left=1.5cm, right=1.5cm, top=1.785cm, bottom=2.0cm]{geometry}
\usepackage{mathptmx}
\usepackage{graphicx}
\usepackage{lastpage}
\usepackage[format=plain,justification=justified,singlelinecheck=false,font={stretch=1.125,small,sf},labelfont=bf,labelsep=space]{caption}
\usepackage{float, booktabs, multirow, array}
\usepackage[T1]{fontenc}
\usepackage{amsmath, lmodern}
\usepackage[detect-all=true,group-minimum-digits=4]{siunitx}
\usepackage{enumitem}
\usepackage{hyperref}
\usepackage{cleveref}
\usepackage{xr-hyper}  % Cross-reference to main document
\usepackage[english]{babel}
\addto{\captionsenglish}{\renewcommand{\refname}{Notes and references}}
\externaldocument[M-]{Article_Digital_Discovery}  % Prefix M- for all main text references

% Prefix SI figures/tables
\renewcommand{\thefigure}{S\arabic{figure}}
\renewcommand{\thetable}{S\arabic{table}}
\renewcommand{\theequation}{S\arabic{equation}}
\setcounter{figure}{0}
\setcounter{table}{0}

\begin{document}

\twocolumn[{
\begin{center}
  {\large\textbf{Supporting Information}}\\[4pt]
  {\large How Far Can Structure-Based Machine Learning Predict Experimental
  Singlet--Triplet Gaps? An Honest Benchmark for TADF Emitter Triage}\\[6pt]
  Jean-Pierre Tchapet Njafa,$^{a}$* Steve Cabrel Teguia Kouam,$^{b}$ Patrick Sorrel Mvoto Kongo,$^{a}$ Panebei Samafou,$^{c}$ and Serge Guy Nana Engo$^{a}$\\[4pt]
  {\small $^a$Department of Physics, Faculty of Science, University of Yaound\'{e}~1, Yaound\'{e}~812, Cameroon.}\\
  {\small $^b$Department of Physics, Faculty of Science, University of Douala, Douala~24157, Cameroon.}\\
  {\small $^c$Department of Medical Imaging and Radiation Sciences, Faculty of Medicine and Health Sciences, Universit\'{e} de Sherbrooke, 3001, 12th Avenue Nord, Sherbrooke, QC J1H~5N4, Canada.}\\[4pt]
  {\small *E-mail: jean-pierre.tchapet@facsciences-uy1.cm; ORCID: https://orcid.org/0000-0002-1936-8353}
\end{center}
\vspace{8pt}
}]

% ─────────────────────────────────────────────────────────────────────────────
\section*{S1. Computational details}
% ─────────────────────────────────────────────────────────────────────────────

\subsection*{S1.1 Semi-empirical calculations}
Ground-state geometries were optimised with GFN2-xTB~\cite{Bannwarth2019}
(\texttt{xtb} 6.5.1, default convergence: $10^{-6}$~Ha SCF). Lowest singlet and
triplet excitations were obtained with sTDA~\cite{Grimme2017}. These
semi-empirical excited states are used only to derive the natural transition
orbital (NTO) descriptors of \cref{tab:descriptors}; they are not used as a
prediction target (the semi-empirical gap is circular with respect to such
features; see main text).

\subsection*{S1.2 DFT reference calculations}
Functional-dependence benchmarks (\cref{tab:dft_benchmark}) used
B3LYP/6-31G(d) (Gaussian~16~\cite{Gaussian16}) and CAM-B3LYP/def2-TZVP
(ORCA~6.1.0~\cite{Neese2020}), in gas phase and toluene (CPCM).

% ─────────────────────────────────────────────────────────────────────────────
\section*{S2. Natural transition orbital descriptors}
% ─────────────────────────────────────────────────────────────────────────────
NTOs~\cite{Martin2003} were obtained by singular-value decomposition of the
transition density matrix; the dominant hole--electron pair
($\lambda_\text{NTO}>0.9$) defines the spatial descriptors. The \num{35}
NTO descriptors---which require the sTDA excited-state calculation but exclude the
$S_1$/$T_1$ energy scalars---are enumerated in \cref{tab:descriptors}. For each of
the lowest singlet ($S_1$) and triplet ($T_1$) states we computed the
hole--electron spatial overlap $S_\text{he}$, the charge-transfer number CT, the
donor/acceptor NTO weights $\lambda_D$/$\lambda_A$, the hole--electron separation
$\Delta r$, and the localisation indices (hole-on-acceptor, particle-on-donor).
State differences ($\Delta X = X_{S_1}-X_{T_1}$) and their magnitudes $|\Delta X|$
were added because the $S_1$--$T_1$ exchange splitting depends on how the spatial
distributions \emph{change} between states. The oscillator strength $f_{S_1}$ and the
log-transformed absorption proxies are the only intensity-related quantities; no
energy scalar enters the feature set.

\begin{table}[htbp]
  \centering
  \caption{The \num{35} NTO spatial descriptors used as the structure-derived
    feature set. $X_{S_1}$ and $X_{T_1}$ denote the quantity evaluated in the
    singlet and triplet state; $\Delta X = X_{S_1}-X_{T_1}$. Energy scalars
    ($E_{S_1}$, $E_{T_1}$, and the gap) are deliberately excluded.}
  \label{tab:descriptors}
  \footnotesize
  \begin{tabular}{p{0.30\columnwidth}p{0.60\columnwidth}}
    \toprule
    Family & Descriptors \\
    \midrule
    Hole--electron overlap (5)
      & $S_\text{he}(S_1)$, $S_\text{he}(T_1)$, $\Delta S_\text{he}$,
        $|\Delta S_\text{he}|$, \texttt{Char\_diff\_squared} \\
    \addlinespace
    Charge-transfer number (4)
      & $\text{CT}(S_1)$, $\text{CT}(T_1)$, $\Delta\text{CT}$, $|\Delta\text{CT}|$ \\
    \addlinespace
    Donor weight (4)
      & $\lambda_D(S_1)$, $\lambda_D(T_1)$, $\Delta\lambda_D$, $|\Delta\lambda_D|$ \\
    \addlinespace
    Acceptor weight (4)
      & $\lambda_A(S_1)$, $\lambda_A(T_1)$, $\Delta\lambda_A$, $|\Delta\lambda_A|$ \\
    \addlinespace
    CT separation (4)
      & $\Delta r(S_1)$, $\Delta r(T_1)$, $\Delta(\Delta r)$, $|\Delta(\Delta r)|$ \\
    \addlinespace
    Localisation (4)
      & hole-on-A$(S_1)$, particle-on-D$(S_1)$,
        hole-on-A$(T_1)$, particle-on-D$(T_1)$ \\
    \addlinespace
    NTO overlap composites (7)
      & $S_\text{NTO}^{\text{sum}}$, $S_\text{NTO}^{\text{prod}}$,
        $S_\text{NTO}^{\text{ratio}}$, $\Delta S_\text{NTO}$,
        $|\Delta S_\text{NTO}|$, NTO overlap$(S_1)$, NTO overlap$(T_1)$ \\
    \addlinespace
    Intensity proxies (3)
      & $f_{S_1}$, $\log|A_{S_1}|$, $\log|A_{T_1}|$ \\
    \bottomrule
  \end{tabular}
\end{table}

% ─────────────────────────────────────────────────────────────────────────────
\section*{S3. Machine learning details}
% ─────────────────────────────────────────────────────────────────────────────

\subsection*{S3.1 Target and dataset}
The target is the experimental $\Delta E_\text{ST}$, extracted from the
literature and matched by compound name; multiple reports were reduced to the
median. The benchmark contains \num{231} molecules spanning \num{212}
Bemis--Murcko scaffolds~\cite{Bemis1996}, so scaffold-based splits place almost
every molecule on a unique chemotype.

The corpus was assembled by a name-based join, which is a fragile step; we therefore
state the rule and its accounting in full (\cref{tab:provenance}). Starting from
\num{1520} extracted literature rows, we kept rows whose property specifier
contained ``EST'' (\num{1490}), parsed the reported value (all \num{1490} parsed),
and retained values in the physical range $[-0.3, 1.5]$~eV (all \num{1490}). Each
row's compound-name list was exploded and matched by exact string equality to the
identifiers of molecules for which we had computed NTO descriptors, giving
\num{363} name--value matches over \num{231} unique molecules. Of these,
\num{102} molecules carried more than one literature report; we collapsed each to
its median. The inter-report spread (standard deviation across reports for a given
molecule) had a median of \qty{0}{\electronvolt} but reached
\qty{0.564}{\electronvolt} in the worst case, which bounds the achievable accuracy
from below. Every matched molecule had an RDKit-parseable SMILES (zero dropped),
yielding the final set of \num{231} molecules over \num{212} scaffolds. The script
\texttt{code/si\_supplementary\_data.py} regenerates this table.

\begin{table}[htbp]
  \centering
  \caption{Provenance of the experimental $\Delta E_\text{ST}$ corpus: rows
    retained or dropped at each stage of the name-based join.}
  \label{tab:provenance}
  \footnotesize
  \begin{tabular}{lc}
    \toprule
    Stage & Count \\
    \midrule
    Extracted literature rows                 & 1520 \\
    Specifier contains ``EST''                & 1490 \\
    Value parsed and finite                   & 1490 \\
    Value in physical range $[-0.3,1.5]$~eV   & 1490 \\
    Name--value matches to descriptor set     & 363 \\
    Unique molecules matched                  & 231 \\
    \quad of which with $>1$ report           & 102 \\
    Dropped: no RDKit-parseable SMILES        & 0 \\
    \midrule
    \textbf{Final benchmark molecules}        & \textbf{231} \\
    \textbf{Final Bemis--Murcko scaffolds}    & \textbf{212} \\
    \bottomrule
  \end{tabular}
\end{table}

The corpus records no measurement medium: the solvent and atmosphere fields are
empty for all \num{1490} EST rows, so the target cannot be stratified by solvent.
The label heterogeneity is instead quantified directly from the spread among
independent reports of the same molecule---the empirical footprint of the
uncontrolled solvent, host and method effects that dominate experimental
$\Delta E_\text{ST}$~\cite{Olivier2017} (\cref{tab:solvent}). Across the \num{102}
molecules with repeat reports, the inter-report range exceeds the model MAE
(\qty{0.096}{\electronvolt}) for \qty{19.6}{\percent} of molecules, and a single
randomly chosen report sits \qty{0.072}{\electronvolt} (RMS) from its molecule's
median. This label noise is comparable to the model error, so it---not model
capacity---sets the accuracy ceiling. Regenerated by
\texttt{code/solvent\_analysis.py}.

\begin{table}[htbp]
  \centering
  \caption{Label heterogeneity of the experimental target, measured from
    within-molecule disagreement among independent literature reports
    (\num{102} molecules with $>1$ report). Range is max$-$min; SD is the standard
    deviation across a molecule's reports (the latter matches \cref{tab:provenance}).}
  \label{tab:solvent}
  \footnotesize
  \begin{tabular}{lc}
    \toprule
    Quantity (eV unless noted) & Value \\
    \midrule
    Molecules with $>1$ report                 & 102 \\
    Inter-report range: mean / median          & 0.062 / 0.000 \\
    Inter-report range: 90th pct.\ / max       & 0.17 / 1.09 \\
    Fraction with range $>$ model MAE          & 19.6\% \\
    Inter-report SD: mean / max                & 0.036 / 0.564 \\
    Single report vs median (RMS)              & 0.072 \\
    \midrule
    Model MAE (reference)                      & 0.096 \\
    \bottomrule
  \end{tabular}
\end{table}

\subsection*{S3.2 Models and evaluation}
\emph{Model.} Random forests~\cite{Breiman2001} (\num{400} trees,
scikit-learn~\cite{Pedregosa2011}) were compared against ridge regression,
support-vector regression and a mean-value baseline; the random forest was best on
every feature set. The forest used \num{400} estimators with no depth cap, $\sqrt{p}$
features per split, and a fixed seed; these defaults were not tuned per fold, so the
reported scores carry no test-set selection bias.

\emph{Cross-validation.} Performance is reported under Bemis--Murcko scaffold
\texttt{GroupKFold} ($k=5$): each fold holds out whole scaffolds, so no scaffold
appears in both training and test, which removes the optimistic leakage of a
random split on a set where \num{212} scaffolds cover \num{231} molecules.

\emph{Intervals and importance.} \qty{95}{\percent} confidence intervals come from
\num{2000} bootstrap resamples of the pooled out-of-fold predictions
(\cref{tab:model_metrics}). Feature importance used exact TreeSHAP~\cite{Lundberg2017}.

\begin{table}[htbp]
  \centering
  \caption{Prediction of experimental $\Delta E_\text{ST}$ (random forest,
    scaffold cross-validation, \num{231} molecules); $R^2$ with \qty{95}{\percent}
    bootstrap CI. The 2D-structure-only sets (no quantum chemistry) match the NTO
    set (which needs an sTDA run); structure-only $R^2$ intervals include zero.}
  \label{tab:model_metrics}
  \small
  \begin{tabular}{lccc}
    \toprule
    Feature set & MAE (eV) & $R^2$ (95\% CI) & $\rho$ \\
    \midrule
    Morgan FP (2D structure only)   & 0.091 & 0.27 ($-0.19$, 0.53) & 0.31 \\
    RDKit desc.\ (2D structure only) & 0.100 & 0.31 ($-0.11$, 0.51) & 0.26 \\
    NTO spatial (needs sTDA)        & 0.096 & 0.25 (0.03, 0.45) & 0.36 \\
    Energies + NTO (reference)      & 0.097 & 0.24 & 0.34 \\
    Mean-value baseline             & 0.107 & 0.00 & --- \\
    Raw sTDA gap                    & 0.251 & $-2.22$ & 0.24 \\
    \bottomrule
  \end{tabular}
\end{table}

\Cref{tab:metrics_ci} collects all accuracy and association metrics for the
headline NTO model with their \qty{95}{\percent} bootstrap confidence intervals,
together with the power analysis for the rank-correlation test. Both Spearman
$\rho$ and Pearson $r$ exclude zero, while the $R^2$ interval does not; the
test for $\rho = 0.36$ is fully powered at $n = 231$ (only $n \approx 60$ is
required for \qty{80}{\percent} power), so the wide $R^2$ interval is a property
of the variance-explained statistic under outlier influence rather than a
sample-size limitation. Values regenerated by
\texttt{code/r2\_ci\_power\_analysis.py}.

\begin{table}[htbp]
  \centering
  \caption{Headline NTO random-forest model: accuracy and association metrics
    with \qty{95}{\percent} percentile-bootstrap CIs (\num{5000} resamples of the
    out-of-fold predictions, \num{231} molecules), and a Fisher-$z$ power analysis
    for the correlation test (two-sided, $\alpha = 0.05$).}
  \label{tab:metrics_ci}
  \small
  \begin{tabular}{lcc}
    \toprule
    Metric & Point estimate & 95\% CI \\
    \midrule
    MAE (eV)            & 0.096 & (0.082, 0.109) \\
    $R^2$               & 0.25  & (0.03, 0.45) \\
    Spearman $\rho$     & 0.36  & (0.24, 0.47) \\
    Pearson $r$         & 0.50  & (0.26, 0.68) \\
    \midrule
    \multicolumn{3}{l}{\emph{Power analysis (correlation test, $\rho = 0.36$)}} \\
    Spearman $\rho$ $p$-value     & \multicolumn{2}{c}{$\num{2.6e-8}$} \\
    Achieved power at $n = 231$   & \multicolumn{2}{c}{$\approx 1.0$} \\
    $n$ for \qty{80}{\percent} power & \multicolumn{2}{c}{60} \\
    $n$ for \qty{90}{\percent} power & \multicolumn{2}{c}{79} \\
    \bottomrule
  \end{tabular}
\end{table}

A scaffold-CV learning curve (\cref{fig:learning}) tests whether the
\num{231}-molecule set is simply too small. The cross-validated MAE shows no
systematic decrease beyond $n \approx 90$: it varies by \qty{5}{\percent} across
$n = 46$--\num{231}, within the fold-to-fold standard deviation
($\approx \qty{0.013}{\electronvolt}$), while the training MAE stays near
\qty{0.035}{\electronvolt}. The persistent train--CV gap that does \emph{not} close
with more data is the signature of an irreducible label-noise floor rather than an
undersampled model; within this data regime, additional molecules of the same
provenance would not lower the error. Regenerated by
\texttt{code/learning\_curve\_analysis.py}.

\begin{figure}[htbp]
  \centering
  \includegraphics[width=0.72\linewidth]{figures/learning_curve.pdf}
  \caption{Scaffold cross-validation learning curve for the headline NTO random
    forest. Shaded band: $\pm 1$ fold standard deviation of the CV-MAE. The CV-MAE
    is flat beyond $n \approx 90$ and remains above the mean-value baseline gap,
    indicating a label-noise ceiling rather than a sample-size limit.}
  \label{fig:learning}
\end{figure}

\subsection*{S3.3 Per-feature SHAP attributions}
\Cref{tab:shap} lists the mean absolute TreeSHAP value of every NTO descriptor for
the random forest trained on the \num{231}-molecule set. Attribution concentrates
on a handful of descriptors: the singlet hole--electron overlap $S_\text{he}(S_1)$
alone accounts for \qty{29}{\percent} of the total attribution, and the top five
descriptors ($S_\text{he}(S_1)$, $S_\text{NTO}^{\text{ratio}}$, $f_{S_1}$,
$\Delta r(S_1)$, $S_\text{he}(T_1)$) for \qty{60}{\percent}. The \num{12}
lowest-ranked descriptors together contribute under \qty{3}{\percent}; the two
raw NTO-overlap and the squared-character features are effectively unused. This is
consistent with the main-text finding that the gap is governed by a small number
of donor--acceptor overlap quantities rather than by the full \num{35}-dimensional
descriptor vector.

\begin{table}[htbp]
  \centering
  \caption{Per-feature importance (mean $|\text{SHAP}|$, in meV) of the \num{35}
    NTO descriptors for the random forest on the \num{231}-molecule experimental
    set, ranked high to low. Regenerated by
    \texttt{code/si\_supplementary\_data.py}.}
  \label{tab:shap}
  \footnotesize
  \begin{tabular}{lr@{\hskip 2em}lr}
    \toprule
    Descriptor & meV & Descriptor & meV \\
    \midrule
    $S_\text{he}(S_1)$        & 36.9 & $|\Delta\text{CT}|$            & 1.6 \\
    $S_\text{NTO}^{\text{ratio}}$ & 14.5 & hole-on-A$(T_1)$          & 1.5 \\
    $f_{S_1}$                 & 11.5 & $\Delta r(T_1)$               & 1.5 \\
    $\Delta r(S_1)$           &  7.9 & $\text{CT}(S_1)$              & 1.4 \\
    $S_\text{he}(T_1)$        &  6.9 & $|\Delta\lambda_A|$          & 1.3 \\
    $\log|A_{T_1}|$           &  6.9 & $\Delta\text{CT}$            & 1.2 \\
    $\Delta(\Delta r)$        &  4.8 & particle-on-D$(S_1)$         & 0.7 \\
    $\Delta S_\text{he}$      &  4.3 & $\lambda_D(S_1)$             & 0.7 \\
    $|\Delta(\Delta r)|$      &  4.1 & particle-on-D$(T_1)$         & 0.6 \\
    $|\Delta S_\text{he}|$    &  4.0 & $\lambda_A(T_1)$             & 0.6 \\
    $|\Delta\lambda_D|$       &  3.1 & $\lambda_A(S_1)$             & 0.5 \\
    $\log|A_{S_1}|$           &  3.0 & hole-on-A$(S_1)$             & 0.5 \\
    $\Delta\lambda_D$         &  2.6 & $\Delta S_\text{NTO}$        & 0.1 \\
    $\Delta\lambda_A$         &  2.0 & $S_\text{NTO}^{\text{sum}}$  & 0.1 \\
    $\lambda_D(T_1)$          &  2.0 & NTO overlap$(T_1)$           & 0.1 \\
    $\text{CT}(T_1)$          &  1.7 & NTO overlap$(S_1)$           & 0.0 \\
    \multicolumn{2}{l}{} & $|\Delta S_\text{NTO}|$         & 0.0 \\
    \multicolumn{2}{l}{} & $S_\text{NTO}^{\text{prod}}$    & 0.0 \\
    \multicolumn{2}{l}{} & \texttt{Char\_diff\_squared}    & 0.0 \\
    \bottomrule
  \end{tabular}
\end{table}

The interpretive value of the physics-informed descriptors is visible when their
attribution is placed beside that of the equally accurate Morgan-fingerprint model
(\cref{fig:shapcompare}). The ten highest-attribution NTO descriptors carry
\qty{79}{\percent} of the total and are all named physical quantities; the ten
highest-attribution Morgan bits carry only \qty{30}{\percent}, spread over
\num{2048} anonymous hashed substructure bits that resist mechanistic reading.
Equal predictive accuracy thus buys interpretability only in the NTO
representation. Regenerated by \texttt{code/shap\_morgan\_vs\_nto.py}.

\begin{figure}[htbp]
  \centering
  \includegraphics[width=\linewidth]{figures/shap_morgan_vs_nto.pdf}
  \caption{Top-ten SHAP attributions for the two equally accurate models. (A) NTO
    random forest: attribution concentrates on named donor--acceptor overlap and
    charge-transfer descriptors. (B) Morgan random forest: attribution is diffuse
    across anonymous fingerprint bits (\texttt{fp}$n$), which carry no direct
    chemical meaning.}
  \label{fig:shapcompare}
\end{figure}

To test whether this attribution pattern is a stable signal rather than an
artefact of one particular fit, we refit the random forest on \num{200} bootstrap
resamples of the \num{231}-molecule set and recorded the top-five descriptors of
each (\cref{tab:shapboot}). The leading descriptor $S_\text{he}(S_1)$ appears in
the top five of \qty{99.5}{\percent} of resamples; on average \num{3.7} of the five
canonical descriptors recur per resample, and \qty{64.5}{\percent} of resamples
recover at least four of the five. The attribution is therefore reproducible at the
top of the ranking, even though the model's ranking power on the gap itself is
modest; we accordingly read these attributions as mechanistic hypotheses for
experimental test, not validated causal claims. Regenerated by
\texttt{code/shap\_bootstrap\_stability.py}.

\begin{table}[htbp]
  \centering
  \caption{Bootstrap stability of the top-five SHAP descriptors: fraction of
    \num{200} resampled-and-refit models in which each canonical descriptor remains
    in the top five.}
  \label{tab:shapboot}
  \footnotesize
  \begin{tabular}{lc}
    \toprule
    Descriptor & Top-5 frequency \\
    \midrule
    $S_\text{he}(S_1)$            & 99.5\% \\
    $S_\text{NTO}^{\text{ratio}}$ & 74.0\% \\
    $f_{S_1}$                     & 72.0\% \\
    $\Delta r(S_1)$               & 68.0\% \\
    $S_\text{he}(T_1)$            & 59.5\% \\
    \midrule
    Mean top-5 overlap            & 3.7 / 5 \\
    Resamples recovering $\ge 4$  & 64.5\% \\
    \bottomrule
  \end{tabular}
\end{table}

\subsection*{S3.4 Scaffold-split characterisation}
Because \num{212} Bemis--Murcko scaffolds cover \num{231} molecules,
\qty{92.9}{\percent} of scaffolds are singletons, and a reviewer may reasonably ask
whether scaffold \texttt{GroupKFold} differs at all from a random split. It does, in
the way that matters for leakage: by construction no scaffold appears in both
training and test (zero shared groups in every fold), whereas a random split of the
same data leaks \numrange{3}{8} scaffolds per fold across the train/test boundary
(\cref{tab:scaffold}). The numerical effect on difficulty is nonetheless small---the
mean nearest-neighbour Tanimoto similarity between a test molecule and its closest
training molecule falls only from \num{0.615} (random) to \num{0.592} (scaffold)---so
we do not claim the split makes the task dramatically harder; we claim it removes a
specific leakage channel.

That the model learns chemical signal rather than scaffold-group artefacts is
established by a label-permutation test: the target was shuffled \num{1000} times and
the scaffold-CV MAE recomputed under each permutation. The true MAE
(\qty{0.0956}{\electronvolt}) lies far in the left tail of the null distribution
(mean \qty{0.121}{\electronvolt}, SD \qty{0.0043}{\electronvolt}, \qty{95}{\percent}
interval \numrange{0.112}{0.129}), giving $p = 0.001$. Regenerated by
\texttt{code/scaffold\_split\_analysis.py}.

\begin{table}[htbp]
  \centering
  \caption{Scaffold \texttt{GroupKFold} versus a random 5-fold split on the same
    \num{231} molecules. Scaffold splitting removes train/test scaffold overlap; the
    label-permutation test confirms the fitted model is far from the no-signal null.}
  \label{tab:scaffold}
  \footnotesize
  \begin{tabular}{lcc}
    \toprule
    Quantity & Scaffold split & Random split \\
    \midrule
    Shared scaffolds per fold        & 0 (all folds) & 3--8 \\
    Test--train NN Tanimoto (mean)   & 0.592 & 0.615 \\
    \midrule
    \multicolumn{3}{l}{\emph{Label-permutation test (1000 shuffles, scaffold CV)}} \\
    True CV-MAE (eV)                 & \multicolumn{2}{c}{0.0956} \\
    Null CV-MAE: mean (95\% CI) (eV) & \multicolumn{2}{c}{0.121 (0.112, 0.129)} \\
    $p$-value (left tail)            & \multicolumn{2}{c}{0.001} \\
    \bottomrule
  \end{tabular}
\end{table}

\subsection*{S3.5 Reproducibility}
The scripts \texttt{code/finalize\_model.py} and
\texttt{code/rebuild\_structural\_model.py} regenerate \cref{tab:model_metrics}
and Fig.~1 of the main text from the deposited data.

% ─────────────────────────────────────────────────────────────────────────────
\section*{S4. DFT level benchmark}
% ─────────────────────────────────────────────────────────────────────────────
Across \num{7} representative emitters in gas and toluene, B3LYP/6-31G(d) and
CAM-B3LYP/def2-TZVP differ in $\Delta E_\text{ST}$ by a mean absolute
\qty{0.30}{\electronvolt} (up to \qty{0.58}{\electronvolt}), with sign reversals
of the relative ordering (e.g.\ DMAC-DPS). This functional dependence sets a
physical ceiling on the meaning of any sub-\qty{0.1}{\electronvolt} accuracy
claim (\cref{tab:dft_benchmark}).

\begin{table}[htbp]
  \centering
  \caption{Functional dependence of $\Delta E_\text{ST}$ (eV): B3LYP/6-31G(d)
    vs CAM-B3LYP/def2-TZVP.}
  \label{tab:dft_benchmark}
  \small
  \begin{tabular}{llccc}
    \toprule
    Molecule & Phase & B3LYP & CAM-B3LYP & $\Delta$ \\
    \midrule
    4CzIPN & gas & 0.192 & 0.361 & +0.169 \\
    4CzIPN & toluene & 0.206 & 0.335 & +0.129 \\
    ACRFLCN & gas & 0.007 & 0.590 & +0.583 \\
    ACRFLCN & toluene & 0.000 & 0.548 & +0.548 \\
    ACRSA & gas & 0.034 & 0.447 & +0.413 \\
    ACRSA & toluene & 0.000 & 0.372 & +0.372 \\
    BACN & gas & 0.524 & 0.991 & +0.467 \\
    BACN & toluene & 0.521 & 0.921 & +0.400 \\
    DMAC-DPS & gas & 0.489 & 0.188 & $-0.301$ \\
    DMAC-DPS & toluene & 0.444 & 0.207 & $-0.237$ \\
    DMAC-TRZ & gas & 0.098 & 0.261 & +0.163 \\
    DMAC-TRZ & toluene & 0.085 & 0.271 & +0.186 \\
    PXZ-NAI & gas & 0.245 & 0.355 & +0.110 \\
    PXZ-NAI & toluene & 0.289 & 0.360 & +0.071 \\
    \bottomrule
  \end{tabular}
\end{table}

A higher-tier check confirms this functional dependence is not an artefact of the
disparate B3LYP/CAM-B3LYP pair. Across six range-separated and hybrid functionals
applied to \num{27} charge-transfer emitters~\cite{Froitzheim2024}, the singlet--triplet gap \emph{estimate}
($E(\text{S}_1^\text{ROKS})-E(\text{S}_1^\text{UKS})\approx\Delta E_\text{ST}/2$)
disagrees per molecule by a median of \qty{0.049}{\electronvolt} (\num{75}th percentile
\qty{0.10}{\electronvolt}, one outlier excluded). Even among modern functionals the
estimate varies at the \qtyrange{0.05}{0.10}{\electronvolt} level; no absolute gap or
experimental agreement is claimed from these data. Extracted by
\texttt{code/parse\_roks\_benchmark.py} from \texttt{data/theoexp\_roks\_energies.csv}.

The observed error floor is consistent with a root-sum-square combination of three
\emph{independently measured} noise sources (\cref{tab:ceiling}). Treating them as
uncorrelated gives an expected MAE floor of \qty{0.128}{\electronvolt}, bracketing the
observed random-forest MAE of \qty{0.096}{\electronvolt}
(\qty{95}{\percent} CI \numrange{0.082}{0.109}) and confirming the ceiling is set by the
data, not the learner (\cref{fig:ceiling}). Computed by
\texttt{code/ceiling\_decomposition.py}.

\begin{table}[htbp]
  \centering
  \caption{Quantitative decomposition of the accuracy ceiling into three measured
    noise components. The root-sum-square floor is compared with the observed
    random-forest MAE and the mean-value baseline.}
  \label{tab:ceiling}
  \small
  \begin{tabular}{llc}
    \toprule
    & Source & Magnitude (eV) \\
    \midrule
    A & Inter-laboratory report scatter (RMS) & 0.072 \\
    B & Vertical$\rightarrow$adiabatic offset (MAE) & 0.195 \\
    C & Functional dependence (mean spread) & 0.296 \\
    \midrule
    & Root-sum-square floor & 0.128 \\
    & Observed RF MAE & 0.096 \\
    & Mean-value baseline & 0.107 \\
    \bottomrule
  \end{tabular}
\end{table}

\begin{figure}[htbp]
  \centering
  \includegraphics[width=0.72\columnwidth]{figures/ceiling_decomposition.pdf}
  \caption{Accuracy-ceiling decomposition: the three measured noise components (A, B, C)
    and their root-sum-square floor against the observed random-forest MAE.}
  \label{fig:ceiling}
\end{figure}

\subsection*{S4.1 Vertical versus adiabatic gap}
The NTO descriptors are evaluated at the ground-state (vertical) geometry. To bound
the error this introduces relative to a relaxed (adiabatic) treatment, we recomputed
$\Delta E_\text{ST}$ for \num{14} representative emitters in gas and toluene at
CAM-B3LYP/def2-TZVP (D3BJ, RIJCOSX) with excited-state geometry optimisation in
ORCA, and compared with the vertical sTDA-xTB values (\cref{tab:adiabatic}). Over the
\num{27} converged points the vertical gap deviates from the adiabatic reference by a
mean absolute \qty{0.195}{\electronvolt} (RMSE \qty{0.245}{\electronvolt}, worst case
\qty{0.548}{\electronvolt} for 2CzTPE), but the deviation is a systematic offset, not
random scatter: a linear fit of adiabatic on vertical gives slope \num{1.00} and
intercept \qty{0.07}{\electronvolt} ($R^2 = 0.70$), and the rank correlation is
preserved (Spearman $\rho = 0.71$). Geometry relaxation therefore shifts the absolute
scale while leaving the ordering---on which the triage use depends---largely intact.
Regenerated by \texttt{code/adiabatic\_validation\_analysis.py} from
\texttt{data/vertical\_vs\_adiabatic\_comparison.csv}.

\begin{table}[htbp]
  \centering
  \caption{Vertical (sTDA-xTB, ground-state geometry) versus adiabatic
    (CAM-B3LYP/def2-TZVP, excited-state optimisation) $\Delta E_\text{ST}$ (eV) for
    \num{14} emitters. $\Delta = $ adiabatic $-$ vertical. Summary: MAE
    \qty{0.195}{\electronvolt}, slope \num{1.00}, intercept \qty{0.07}{\electronvolt},
    $R^2 = 0.70$, $\rho = 0.71$.}
  \label{tab:adiabatic}
  \footnotesize
  \begin{tabular}{llccc}
    \toprule
    Molecule & Phase & Vertical & Adiabatic & $\Delta$ \\
    \midrule
    4CzIPN & gas & 0.361 & 0.402 & $+0.041$ \\
    4CzIPN & toluene & 0.335 & 0.436 & $+0.101$ \\
    DMAC-TRZ & gas & 0.261 & 0.442 & $+0.181$ \\
    DMAC-TRZ & toluene & 0.271 & 0.638 & $+0.367$ \\
    DMAC-DPS & gas & 0.188 & 0.385 & $+0.197$ \\
    DMAC-DPS & toluene & 0.207 & 0.558 & $+0.351$ \\
    PXZ-NAI & gas & 0.355 & 0.243 & $-0.112$ \\
    PXZ-NAI & toluene & 0.360 & 0.491 & $+0.131$ \\
    TPA-APy & gas & 0.876 & 0.753 & $-0.123$ \\
    TPA-APy & toluene & 0.765 & 0.895 & $+0.130$ \\
    BACN & gas & 0.991 & 1.129 & $+0.138$ \\
    BACN & toluene & 0.921 & 1.063 & $+0.142$ \\
    BMZ-TZ & gas & 0.938 & 0.410 & $-0.528$ \\
    BMZ-TZ & toluene & 0.933 & 0.700 & $-0.233$ \\
    2CzPN & gas & 0.613 & 0.560 & $-0.053$ \\
    2CzPN & toluene & 0.621 & 0.652 & $+0.031$ \\
    APPT-PXZ & gas & 0.350 & 0.331 & $-0.019$ \\
    2PXZP & gas & 0.467 & 0.384 & $-0.083$ \\
    2PXZP & toluene & 0.485 & 0.639 & $+0.154$ \\
    ACRSA & gas & 0.447 & 0.204 & $-0.243$ \\
    ACRSA & toluene & 0.372 & 0.677 & $+0.305$ \\
    ACRFLCN & gas & 0.590 & 0.302 & $-0.288$ \\
    ACRFLCN & toluene & 0.548 & 0.576 & $+0.028$ \\
    2CzTPE & gas & 1.143 & 1.691 & $+0.548$ \\
    2CzTPE & toluene & 1.042 & 1.535 & $+0.493$ \\
    BBPA & gas & 1.521 & 1.641 & $+0.120$ \\
    BBPA & toluene & 1.433 & 1.562 & $+0.129$ \\
    \bottomrule
  \end{tabular}
\end{table}

\subsection*{S4.2 Triplet manifold (T1--T2)}\label{sec:triplet_manifold}
A small $S_1$--$T_1$ gap does not guarantee fast reverse intersystem crossing; the
$T_1$--$T_2$ splitting and the wider triplet manifold matter too~\cite{Keruckiene2024}.
The same vertical CAM-B3LYP/def2-TZVP TD-DFT calculations that supplied the adiabatic
benchmark also resolve ten triplet roots, so we report $T_1$ and $T_2$ directly
(\cref{tab:manifold}). Across the \num{27} converged emitter--phase points the $T_1$--$T_2$
gap ranges from \qty{0.018}{\electronvolt} (4CzIPN) to \qty{1.50}{\electronvolt} (BBPA),
mean \qty{0.41}{\electronvolt}; \num{13} of \num{27} fall below \qty{0.3}{\electronvolt}.
The manifold therefore varies widely \emph{among molecules with similar $\Delta E_\text{ST}$}
(e.g.\ DMAC-DPS and DMAC-TRZ have near-equal $S_1$--$T_1$ gaps but $T_1$--$T_2$ of
\qtyrange{0.06}{0.26}{\electronvolt}), which is exactly the information a scalar
$\Delta E_\text{ST}$ triage cannot encode and a direction where the excited-state pipeline
adds genuine value. Regenerated by \texttt{code/extract\_triplet\_manifold.py} from the
\texttt{orca\_package} outputs.

\begin{table}[htbp]
  \centering
  \caption{Vertical CAM-B3LYP/def2-TZVP TDA energies (eV): $S_1$, $T_1$, $T_2$, the
    singlet--triplet gap $\Delta E_\text{ST}=S_1-T_1$, and the triplet-manifold gap
    $T_1$--$T_2$, for \num{14} emitters in gas and toluene.}
  \label{tab:manifold}
  \scriptsize
  \begin{tabular}{llccccc}
    \toprule
    Molecule & Phase & $S_1$ & $T_1$ & $T_2$ & $\Delta E_\text{ST}$ & $T_1$--$T_2$ \\
    \midrule
    2CzPN & gas & 3.652 & 3.039 & 3.464 & 0.613 & 0.425 \\
    2CzPN & toluene & 3.632 & 3.011 & 3.459 & 0.621 & 0.448 \\
    2CzTPE & gas & 3.496 & 2.353 & 2.767 & 1.143 & 0.414 \\
    2CzTPE & toluene & 3.411 & 2.369 & 2.788 & 1.042 & 0.419 \\
    2PXZP & gas & 3.292 & 2.825 & 3.080 & 0.467 & 0.255 \\
    2PXZP & toluene & 3.343 & 2.858 & 3.139 & 0.485 & 0.281 \\
    4CzIPN & gas & 3.377 & 3.016 & 3.050 & 0.361 & 0.034 \\
    4CzIPN & toluene & 3.343 & 3.008 & 3.026 & 0.335 & 0.018 \\
    ACRFLCN & gas & 3.597 & 3.007 & 3.548 & 0.590 & 0.541 \\
    ACRFLCN & toluene & 3.546 & 2.998 & 3.536 & 0.548 & 0.538 \\
    ACRSA & gas & 3.769 & 3.322 & 3.513 & 0.447 & 0.191 \\
    ACRSA & toluene & 3.746 & 3.374 & 3.500 & 0.372 & 0.126 \\
    APPT-PXZ & gas & 2.701 & 2.351 & 2.734 & 0.350 & 0.383 \\
    BACN & gas & 3.713 & 2.722 & 3.533 & 0.991 & 0.811 \\
    BACN & toluene & 3.629 & 2.708 & 3.503 & 0.921 & 0.795 \\
    BBPA & gas & 3.553 & 2.032 & 3.525 & 1.521 & 1.493 \\
    BBPA & toluene & 3.458 & 2.025 & 3.526 & 1.433 & 1.501 \\
    BMZ-TZ & gas & 3.730 & 2.792 & 3.180 & 0.938 & 0.388 \\
    BMZ-TZ & toluene & 3.896 & 2.963 & 3.142 & 0.933 & 0.179 \\
    DMAC-DPS & gas & 3.637 & 3.449 & 3.508 & 0.188 & 0.059 \\
    DMAC-DPS & toluene & 3.678 & 3.471 & 3.538 & 0.207 & 0.067 \\
    DMAC-TRZ & gas & 3.509 & 3.248 & 3.504 & 0.261 & 0.256 \\
    DMAC-TRZ & toluene & 3.533 & 3.262 & 3.509 & 0.271 & 0.247 \\
    PXZ-NAI & gas & 2.688 & 2.333 & 2.836 & 0.355 & 0.503 \\
    PXZ-NAI & toluene & 2.650 & 2.290 & 2.829 & 0.360 & 0.539 \\
    TPA-APy & gas & 3.555 & 2.679 & 2.823 & 0.876 & 0.144 \\
    TPA-APy & toluene & 3.425 & 2.660 & 2.795 & 0.765 & 0.135 \\
    \bottomrule
  \end{tabular}
\end{table}

To read the manifold across the whole corpus rather than the \num{14}-molecule ORCA
subset, we extracted $T_1$ and $T_2$ for all \num{231} emitters from the cheaper
sTDA-xTB pipeline (\texttt{code/extract\_t1\_t2\_rerun.py}; \num{231}/\num{231}
converged, \texttt{data/t1\_t2\_full.csv}). The full-corpus $T_1$--$T_2$ gap has median
\qty{0.14}{\electronvolt} (mean \qty{0.24}{\electronvolt}), and \qty{74}{\percent} of
emitters fall below the \qty{0.3}{\electronvolt} region where a near-degenerate second
triplet can open competing non-radiative channels (\cref{fig:t1t2dist}). Most
low-gap candidates therefore remain exposed to inefficient reverse intersystem
crossing---a hazard a scalar $S_1$--$T_1$ triage cannot see.

This full-corpus manifold must be read \emph{qualitatively only}. Benchmarking the
sTDA-xTB $T_1$--$T_2$ gap against the CAM-B3LYP reference on the six overlapping
molecules gives a mean absolute deviation of \qty{0.237}{\electronvolt}
(Spearman $\rho = 0.43$, $p = 0.40$; \cref{tab:manifold_val}), so the two methods agree
on neither magnitude nor rank at a useful level. Quantitative ranking on RISC propensity
still requires the higher-level calculation; the sTDA manifold serves only to flag that
the triplet-manifold hazard is widespread, not to order candidates by it.

\begin{table}[htbp]
  \centering
  \caption{sTDA-xTB versus CAM-B3LYP $T_1$--$T_2$ gap (eV) for the six molecules present
    in both sets. Independent methods; the comparison is genuine (fresh GFN2 geometries).
    Summary: MAE \qty{0.237}{\electronvolt}, Spearman $\rho = 0.43$ ($p = 0.40$, $n = 6$).}
  \label{tab:manifold_val}
  \small
  \begin{tabular}{lcc}
    \toprule
    Molecule & sTDA-xTB & CAM-B3LYP \\
    \midrule
    2CzPN    & 0.169 & 0.425 \\
    2PXZP    & 0.004 & 0.255 \\
    4CzIPN   & 0.073 & 0.034 \\
    ACRSA    & 0.006 & 0.191 \\
    DMAC-DPS & 0.011 & 0.059 \\
    PXZ-NAI  & 1.146 & 0.503 \\
    \bottomrule
  \end{tabular}
\end{table}

\begin{figure}[htbp]
  \centering
  \includegraphics[width=0.72\columnwidth]{figures/t1_t2_distribution.pdf}
  \caption{Distribution of the sTDA-xTB $T_1$--$T_2$ gap across all \num{231} emitters.
    \qty{74}{\percent} fall below \qty{0.3}{\electronvolt} (dashed line). Read
    qualitatively only (\cref{tab:manifold_val}).}
  \label{fig:t1t2dist}
\end{figure}

% ─────────────────────────────────────────────────────────────────────────────
\section*{S5. Heuristic library enumeration}
% ─────────────────────────────────────────────────────────────────────────────
An extended library of \num{22194} donor--acceptor fragment combinations was
enumerated from literature-validated donors (carbazole, phenoxazine, acridine)
and acceptors (triazine, benzonitrile, sulfone) under simple constraints
(MW $<\qty{900}{\dalton}$, $\geq 2$ aromatic rings, $\geq 1$ N/O/S). Candidates
were ranked by the rule-based TADF score defined in the main text. These are
synthesis-planning leads only: they are not quantum-validated, and the heuristic
score is not a calibrated $\Delta E_\text{ST}$ predictor. Property distributions
of the enumerated set versus the top-scoring subset are shown in the main text.

To make the section actionable rather than decorative, we additionally rank the full
library with the calibrated structure-only Morgan-FP model and attach a split-conformal
\qty{90}{\percent} prediction interval (half-width \qty{0.18}{\electronvolt}, empirical
coverage \num{0.89}; \cref{tab:shortlist}) to every prediction. The complete top-\num{100}
shortlist is deposited (\texttt{data/shortlist\_top100.csv}). We stress the honest reading:
the predicted gaps for the leading candidates (\qtyrange{0.02}{0.05}{\electronvolt}) sit
well inside their own \qty{\pm0.18}{\electronvolt} intervals, so the list \emph{prioritises}
synthesis order; it does not certify any individual molecule. No device-level quantity is
predicted. Generated by \texttt{code/build\_shortlist.py}.

\begin{table}[htbp]
  \centering
  \caption{Top \num{10} of the deposited \num{100}-molecule shortlist: enumerated-library
    candidates ranked by predicted experimental $\Delta E_\text{ST}$ (structure-only
    Morgan-FP RF) with \qty{90}{\percent} split-conformal prediction intervals (eV).
    SMILES truncated for display.}
  \label{tab:shortlist}
  \scriptsize
  \begin{tabular}{llcc}
    \toprule
    ID & SMILES (truncated) & Pred.\ (eV) & 90\% PI (eV) \\
    \midrule
    MOL\_016680 & \texttt{O=C(C[C@@]1(O)C(=O)N(Cc2ccccc2Cl)c…} & 0.023 & [-0.15, 0.20] \\
    MOL\_016681 & \texttt{O=C(C[C@]1(O)C(=O)N(Cc2ccccc2Cl)c2…} & 0.023 & [-0.15, 0.20] \\
    MOL\_002952 & \texttt{O=c1c2ccccc2c(=O)c2c1ccc1c2[nH]c2c…} & 0.034 & [-0.14, 0.21] \\
    MOL\_000892 & \texttt{Nc1cc2c(=O)c3cc(Cl)cc(Cl)c3[nH]c2c…} & 0.035 & [-0.14, 0.21] \\
    MOL\_000969 & \texttt{O=C1c2ccccc2C(=O)c2c1c(Cl)cc1c(=O)…} & 0.038 & [-0.14, 0.22] \\
    MOL\_002522 & \texttt{Cc1c2ccncc2c(C)c2c1[nH]c1c(Cl)cc(C…} & 0.039 & [-0.14, 0.22] \\
    MOL\_017220 & \texttt{Cc1ccccc1CN1C(=O)[C@](O)(CC(=O)c2c…} & 0.039 & [-0.14, 0.22] \\
    MOL\_017221 & \texttt{Cc1ccccc1CN1C(=O)[C@@](O)(CC(=O)c2…} & 0.039 & [-0.14, 0.22] \\
    MOL\_017226 & \texttt{O=C(C[C@@]1(O)C(=O)N(CCc2ccccc2)c2…} & 0.041 & [-0.14, 0.22] \\
    MOL\_017407 & \texttt{O=C(C[C@]1(O)C(=O)N(Cc2ccccc2)c2cc…} & 0.041 & [-0.14, 0.22] \\
    \bottomrule
  \end{tabular}
\end{table}

% ─────────────────────────────────────────────────────────────────────────────
\section*{S6. Robustness analyses}
% ─────────────────────────────────────────────────────────────────────────────

\subsection*{S6.1 Cluster cross-validation}
Because \qty{92.9}{\percent} of scaffolds are singletons, we built non-singleton groups
by Butina clustering (Morgan Tanimoto distance) and ran cluster-grouped CV across a
cutoff scan (\cref{tab:clustercv}). Even at
the coarsest grouping (\num{70} clusters, largest \num{38} molecules), the MAE
(\qty{0.099}{\electronvolt}) is statistically indistinguishable from scaffold-CV
(\qty{0.096}{\electronvolt}) and random-CV (\qty{0.106}{\electronvolt}). Performance is
therefore \emph{robust to the validation scheme}: the headline number is not an artefact of
singleton scaffolds. Regenerated by \texttt{code/cluster\_cv\_analysis.py}.

\begin{table}[htbp]
  \centering
  \caption{Cluster-grouped CV of the headline NTO RF across Butina distance cutoffs,
    with scaffold and random CV for reference.}
  \label{tab:clustercv}
  \footnotesize
  \begin{tabular}{lcccc}
    \toprule
    Grouping & \#groups & largest & MAE (eV) & $\rho$ \\
    \midrule
    Butina 0.3 & 186 & 7  & 0.100 & 0.35 \\
    Butina 0.4 & 153 & 10 & 0.100 & 0.38 \\
    Butina 0.5 & 111 & 15 & 0.104 & 0.31 \\
    Butina 0.6 & 70  & 38 & 0.099 & 0.40 \\
    \midrule
    Scaffold (headline) & 212 & --- & 0.096 & 0.36 \\
    Random 5-fold       & --- & --- & 0.106 & 0.34 \\
    \bottomrule
  \end{tabular}
\end{table}

\subsection*{S6.2 Dimensionality and model capacity}
Two checks address the charge that a \num{2048}-bit fingerprint over-parameterises a
\num{231}-molecule set. First, shrinking the fingerprint to \num{1024} or \num{512} bits
leaves the MAE and the (negative) lower $R^2$ confidence bound essentially unchanged
(\cref{tab:capacity}, top), so the negative $R^2$ tail is intrinsic to the low
variance-explained on this target, not a width artefact; the lowest-dimensional regularised
model (NTO + ElasticNet) brings the lower bound closest to zero ($-0.03$) but does not lift
it. Second, no higher- or different-capacity learner beats the random forest under the same
scaffold CV (\cref{tab:capacity}, bottom): gradient boosting, an MLP and an SVR are all
worse. Model capacity is not the limiting factor. Restricting to the \num{208} molecules
whose inter-report standard deviation is $\le\qty{0.05}{\electronvolt}$ (single-report
molecules included, SD~$=0$) does not lower the MAE (\qty{0.106}{\electronvolt}), so the
floor is not removed by discarding the noisiest labels either. Regenerated by \texttt{code/dimensionality\_check.py} and
\texttt{code/capacity\_and\_noise\_ceiling.py}.

\begin{table}[htbp]
  \centering
  \caption{Top: fingerprint-width sweep (RF) — MAE and \qty{95}{\percent} bootstrap $R^2$
    CI. Bottom: alternative learners under the same scaffold CV. None beats the RF.}
  \label{tab:capacity}
  \footnotesize
  \begin{tabular}{lcc}
    \toprule
    Setting & MAE (eV) & $R^2$ (95\% CI) \\
    \midrule
    Morgan 512 (RF)   & 0.098 & 0.15 ($-0.44$, 0.46) \\
    Morgan 1024 (RF)  & 0.096 & 0.24 ($-0.25$, 0.51) \\
    Morgan 2048 (RF)  & 0.091 & 0.27 ($-0.19$, 0.53) \\
    NTO + ElasticNet  & 0.102 & 0.10 ($-0.03$, 0.17) \\
    \midrule
    RandomForest (headline) & 0.096 & --- \\
    HistGradientBoosting    & 0.105 & --- \\
    MLP (100, 50)           & 0.169 & --- \\
    SVR (RBF)               & 0.120 & --- \\
    ElasticNet              & 0.102 & --- \\
    \bottomrule
  \end{tabular}
\end{table}

\subsection*{S6.3 SHAP cross-check by permutation importance}
An independent, model-agnostic permutation-importance analysis (held-out \qty{30}{\percent},
\num{50} repeats) ranks the singlet hole--electron overlap $S_\text{he}(S_1)$ \emph{first},
in agreement with TreeSHAP; two of the SHAP top-five recur in the permutation top-five. The
leading attribution is thus corroborated by two methods, while the lower-ranked tail is not
--- consistent with our restriction of mechanistic reading to the dominant descriptor.
Regenerated by \texttt{code/permutation\_importance\_check.py}.

\subsection*{S6.4 Enrichment curve and calibrated operating points}
\Cref{fig:enrichment} reports precision (fraction of selected molecules with experimental
gap $\le\qty{0.10}{\electronvolt}$) against screening budget. Enrichment peaks at
\num{1.37}$\times$ over the \qty{49}{\percent} base rate around the top \num{5}--\num{20}\%
and does not improve at more aggressive cuts (the top \num{1}\% is two molecules). A
split-conformal calibration gives a \qty{90}{\percent} interval half-width of
\qty{0.178}{\electronvolt} (empirical coverage \num{0.892}), so a laboratory can pick an
operating point at a stated confidence. Regenerated by \texttt{code/enrichment\_curve.py}.

\begin{figure}[htbp]
  \centering
  \includegraphics[width=0.7\linewidth]{figures/enrichment_curve.pdf}
  \caption{Triage precision versus fraction of the benchmark screened (log $x$). Dashed line:
    the \qty{49}{\percent} random base rate. Enrichment is real but bounded at
    $\approx 1.3$--$1.4\times$.}
  \label{fig:enrichment}
\end{figure}

\subsection*{S6.5 Multi-criteria triage}
A small $S_1$--$T_1$ gap is necessary but not sufficient for efficient delayed
fluorescence: the emitter must also be optically bright and free of a near-degenerate
second triplet. We therefore composed the predicted gap with two auxiliary structure-level
filters---oscillator strength $f_{S_1}\ge\num{0.01}$ (\num{209}/\num{231} pass;
\texttt{code/oscillator\_strength\_model.py}) and sTDA $T_1$--$T_2 > \qty{0.10}{\electronvolt}$
(\cref{tab:manifold_val})---leaving \num{131} molecules that clear all three
(\texttt{code/multi\_criteria\_triage.py}). Ranking within this pool raises the top-\qty{5}{\percent}
precision for low-gap emitters from \num{0.58} (gap alone) to \num{0.71}, i.e.\ enrichment
\num{1.18}$\rightarrow$\num{1.45}$\times$; the two curves converge by the top \qty{50}{\percent}.
The gain is real but small and rests on only a handful of molecules at the sharpest cut, so we
read the multi-criteria filter not as a precision gain but as a physically motivated way
to discard the \qty{43}{\percent} of low-gap candidates that are dark or manifold-unsafe---a
decision a scalar-gap triage cannot make. The $T_1$--$T_2$ term inherits the qualitative-only
caveat of \cref{tab:manifold_val}.

\subsection*{S6.6 Dataset expansion stress-test}
To test whether the benchmark is limited by corpus size, we resolved and featurised
\num{624} additional donor--acceptor emitters from the experimental literature
(name-to-structure resolution followed by the identical GFN2-xTB/sTDA natural-transition-orbital
pipeline; \texttt{code/expand\_dataset.py}, \texttt{code/featurize\_nto\_local.py}). After
discarding \num{127} structures already present in the training set, \num{497} genuinely new
molecules extend the corpus to \num{728} (\cref{tab:expansion}). Enlarging the set does
\emph{not} improve structure-only prediction: under the same scaffold CV the MAE rises from
\qty{0.098}{\electronvolt} to \qty{0.124}{\electronvolt} and $R^2$ falls to \num{-0.01} (the
curated-set MAE of \qty{0.098}{\electronvolt} reproduces the headline \qty{0.096}{\electronvolt}
within single-seed variation). Restricting the new molecules to those with replicate
measurements (\num{234} with $\ge\num{2}$ independent reports) recovers part of the loss
(MAE \qty{0.110}{\electronvolt}, $R^2$ \num{0.05}), placing one component of the degradation in
label noise from automated literature extraction; the residual gap to the baseline reflects the
greater chemical diversity of the larger set. A threefold expansion therefore does not lift the
accuracy ceiling --- it exposes it. Regenerated by \texttt{code/expansion\_stress\_test.py}.

\begin{table}[htbp]
  \centering
  \caption{Scaffold-CV of the NTO random forest as the corpus is expanded and as
    new-label confidence is tightened. Expansion does not improve accuracy.}
  \label{tab:expansion}
  \footnotesize
  \begin{tabular}{lcccc}
    \toprule
    Corpus & $n$ & MAE (eV) & $R^2$ & $\rho$ \\
    \midrule
    Curated baseline            & 231 & 0.098 & \phantom{$-$}0.18 & 0.34 \\
    + all new                   & 728 & 0.124 & $-0.01$           & 0.16 \\
    + multi-report new ($\ge 2$) & 465 & 0.110 & \phantom{$-$}0.05 & 0.24 \\
    \bottomrule
  \end{tabular}
\end{table}

% ─────────────────────────────────────────────────────────────────────────────
\section*{S7. Data and code availability}
% ─────────────────────────────────────────────────────────────────────────────
All data (experimental corpus, descriptor tables, predictions), trained models,
and analysis/figure scripts are openly available on Zenodo
(DOI: 10.5281/zenodo.17436069) under a CC-BY-4.0 licence. Computations used
\texttt{xtb} 6.7.1, \texttt{xtb4stda} with \texttt{stda}~1.6.3, Multiwfn~3.8,
Gaussian~16, ORCA~6.1, Python~3.12, scikit-learn~1.8, RDKit~2026.03, and SHAP~0.50.

% ─────────────────────────────────────────────────────────────────────────────
\section*{S8. Inverted singlet--triplet emitters (out of scope)}
% ─────────────────────────────────────────────────────────────────────────────
The benchmark is restricted to conventional TADF emitters with a positive gap
($\Delta E_\text{ST} > 0$). The emerging inverted singlet--triplet (INVEST) class,
which violates Hund's rule through double-excitation character~\cite{Veglianti2025},
lies outside the reach of the present pipeline: GFN2-xTB enforces Aufbau ordering, and
sTDA-xTB cannot reproduce the doubly-occupied LUMO character required for sign inversion,
so a negative gap is structurally impossible for the random forest to predict from the
current feature set. Screening the \num{25482}-record literature database
(\texttt{code/ist\_screening.py}) returns \num{66} reports of a reported inverted gap
across \num{61} distinct compound names (\num{25} DOIs), spanning experimental
$\Delta E_\text{ST}$ from \qty{-0.43}{\electronvolt} to \qty{-0.01}{\electronvolt}
(mean \qty{-0.11}{\electronvolt}; \cref{fig:istdist},
\texttt{data/ist\_candidates.csv}). We do not classify their cores here because the
database entries carry no curated SMILES. Of the \num{61}, five also appear in our
\num{231}-molecule
corpus; the random forest assigns a \emph{positive} gap to all five
(\qtyrange{0.115}{0.210}{\electronvolt}), i.e.\ the wrong sign in \num{5}/\num{5} cases.
This is the expected falsification: INVEST screening requires ROKS- or EOM-CCSD-level
descriptors~\cite{Salvi2025} and is an explicit out-of-scope future axis, not a
capability of a structure-only model trained on Aufbau-ordered features.

\begin{figure}[htbp]
  \centering
  \includegraphics[width=0.72\columnwidth]{figures/ist_delta_est_distribution.pdf}
  \caption{Reported experimental $\Delta E_\text{ST}$ for the \num{66} inverted-gap
    records identified in the literature database. All lie below the Hund's-rule
    boundary (dashed line); none is reachable by the Aufbau-ordered feature set.}
  \label{fig:istdist}
\end{figure}

\bibliographystyle{../MCA_submission/rsc}
\bibliography{references_digital_discovery}

\end{document}
